\documentclass[11pt]{article}

\usepackage[margin=1in]{geometry}
\usepackage{amsmath,amssymb,amsthm}
\usepackage{booktabs}
\usepackage{array}
\usepackage{longtable}
\usepackage{enumitem}
\usepackage{xcolor}
\usepackage{hyperref}
\usepackage{titlesec}

\hypersetup{
    colorlinks=true,
    linkcolor=blue,
    urlcolor=blue,
    citecolor=blue
}

\titleformat{\section}
  {\large\bfseries}
  {\thesection}
  {0.5em}
  {}

\titleformat{\subsection}
  {\normalsize\bfseries}
  {\thesubsection}
  {0.5em}
  {}

\setlist[itemize]{noitemsep, topsep=2pt}
\setlist[enumerate]{noitemsep, topsep=2pt}

\title{
    \textbf{EN.553.439/639 -- Time Series Analysis}\\
    \large Comprehensive Course Structure and Learning Breakdown
}

\author{Course Structure Derived from the Spring 2026 Syllabus}
\date{}

\begin{document}

\maketitle

\section{Course Overview}

\textbf{Course:} EN.553.439/639 -- Time Series Analysis

\textbf{Credits:} 3

\textbf{Primary Text:}
J. D. Cryer and K.-S. Chan,
\textit{Time Series Analysis with Applications in R},
2nd edition, Springer.

\subsection{Course Description}

The course develops statistical methods for analyzing observations collected
over time. The emphasis is on both time-domain and frequency-domain
approaches, including descriptive techniques, regression analysis,
trend modeling, smoothing, prediction, linear systems, serial correlation,
stationary processes, and spectral analysis.

The principal modeling framework progresses from basic stochastic
processes to linear models, ARMA and ARIMA models, model identification,
parameter estimation, diagnostics, forecasting, seasonal models,
regression/intervention methods, and conditional heteroskedasticity.

\section{Course Goals}

By the end of the course, a student should be able to:

\begin{enumerate}
    \item Describe the statistical structure of a time series.
    \item Quantify serial dependence using covariance, autocovariance,
          and autocorrelation functions.
    \item Distinguish stationary from nonstationary time series.
    \item Identify deterministic and stochastic trends.
    \item Construct and analyze linear time-series models.
    \item Understand and apply AR, MA, and ARMA models.
    \item Transform nonstationary series using differencing and
          variance-stabilizing transformations.
    \item Construct ARIMA models for nonstationary time series.
    \item Identify plausible models using ACF, PACF, and EACF.
    \item Estimate model parameters using Yule--Walker equations,
          maximum likelihood, least squares, and conditional
          least squares.
    \item Evaluate model adequacy through residual diagnostics.
    \item Produce forecasts from ARIMA-type models.
    \item Construct and interpret prediction intervals.
    \item Model seasonal time series using SARIMA models.
    \item Analyze interventions and outliers in time-series data.
    \item Recognize and diagnose spurious correlation.
    \item Apply pre-whitening when appropriate.
    \item Understand conditional heteroskedasticity and volatility
          models such as ARCH and GARCH.
    \item Develop familiarity with nonlinear and frequency-domain
          methods when time permits.
\end{enumerate}

\section{Prerequisite Knowledge}

The official prerequisite structure requires an appropriate prior
statistics/data-analysis course together with a qualifying probability
or statistics course.

For successful completion of this course, students should be comfortable
with the following mathematical and statistical concepts.

\subsection{Probability}

\begin{itemize}
    \item Random variables and random vectors
    \item Expectation
    \item Variance and covariance
    \item Conditional expectation
    \item Probability distributions
    \item Independence and dependence
\end{itemize}

\subsection{Statistics}

\begin{itemize}
    \item Point estimation
    \item Confidence intervals
    \item Hypothesis testing
    \item Regression analysis
    \item Least-squares estimation
    \item Maximum likelihood
\end{itemize}

\subsection{Mathematics}

\begin{itemize}
    \item Differential calculus
    \item Optimization
    \item Basic integration
    \item Matrix algebra
    \item Solving systems of linear equations
    \item Eigenvalues and eigenvectors
\end{itemize}

\subsection{Computational Skills}

Students should be able to:

\begin{itemize}
    \item Manipulate numerical data
    \item Produce basic statistical plots
    \item Fit statistical models
    \item Interpret statistical output
    \item Work with statistical software, particularly R
\end{itemize}

\section{Course Dependency Structure}

The course can be viewed as the following dependency chain:

\[
\boxed{
\text{Probability/Statistics}
\rightarrow
\text{Time-Series Fundamentals}
\rightarrow
\text{Stationarity}
\rightarrow
\text{ARMA}
\rightarrow
\text{ARIMA}
}
\]

\[
\boxed{
\text{ARIMA}
\rightarrow
\text{Identification}
\rightarrow
\text{Estimation}
\rightarrow
\text{Diagnostics}
\rightarrow
\text{Forecasting}
}
\]

\[
\boxed{
\text{ARIMA}
\rightarrow
\text{Seasonality}
\rightarrow
\text{SARIMA}
\rightarrow
\text{Seasonal Forecasting}
}
\]

\[
\boxed{
\text{Regression}
+
\text{Time-Series Dependence}
\rightarrow
\text{Intervention Analysis}
+
\text{Pre-whitening}
}
\]

\[
\boxed{
\text{Conditional Variance}
\rightarrow
\text{ARCH}
\rightarrow
\text{GARCH}
}
\]

\section{Module 1: Fundamental Tools for Time-Series Analysis}

\subsection{Purpose}

Establish the mathematical language required to describe dependence
between observations occurring at different times.

\subsection{Major Topics}

\begin{itemize}
    \item Mean function
    \item Variance function
    \item Covariance
    \item Autocovariance
    \item Autocorrelation
    \item Sample autocorrelation
\end{itemize}

\subsection{Key Concepts}

For a time series $\{X_t\}$, define

\[
\mu_t = E[X_t],
\]

\[
\operatorname{Var}(X_t)
=
E[(X_t-\mu_t)^2],
\]

and

\[
\gamma(s,t)
=
\operatorname{Cov}(X_s,X_t).
\]

For a stationary process, autocovariance is commonly written as

\[
\gamma(h)
=
\operatorname{Cov}(X_t,X_{t+h}),
\]

and autocorrelation is

\[
\rho(h)
=
\frac{\gamma(h)}{\gamma(0)}.
\]

\subsection{Learning Objectives}

Students should be able to:

\begin{enumerate}
    \item Calculate means and variances of time-series observations.
    \item Explain serial dependence.
    \item Calculate covariance and autocovariance.
    \item Interpret the ACF.
    \item Distinguish ordinary correlation from serial correlation.
    \item Explain how dependence changes with lag.
\end{enumerate}

\subsection{Difficulty}

\textbf{Level: Moderate}

The computational formulas are relatively straightforward, but interpretation
of dependence across lags becomes important for subsequent modeling.

\section{Module 2: Basic Time-Series Models and Stationarity}

\subsection{Major Topics}

\begin{itemize}
    \item Basic stochastic time-series models
    \item Strict stationarity
    \item Weak stationarity
    \item Dependence structure
    \item White noise
\end{itemize}

\subsection{Key Concepts}

A weakly stationary process satisfies

\[
E[X_t] = \mu,
\]

\[
\operatorname{Var}(X_t) = \sigma^2,
\]

and

\[
\operatorname{Cov}(X_t,X_{t+h}) = \gamma(h),
\]

where the covariance depends only on the lag $h$ and not on the
specific time $t$.

Strict stationarity concerns invariance of the complete joint
distribution under time shifts.

\subsection{Learning Objectives}

Students should be able to:

\begin{enumerate}
    \item Define strict stationarity.
    \item Define weak stationarity.
    \item Explain the distinction between strict and weak stationarity.
    \item Determine whether a model satisfies stationarity conditions.
    \item Explain why stationarity is central to ARMA/ARIMA modeling.
    \item Recognize white-noise processes.
\end{enumerate}

\subsection{Difficult Concepts}

\begin{itemize}
    \item Difference between distributional and second-order stationarity
    \item Why dependence can exist while the mean remains constant
    \item Stationarity conditions for stochastic models
\end{itemize}

\textbf{Difficulty: High}

\section{Module 3: Trends and Transformations}

\subsection{Major Topics}

\begin{itemize}
    \item Deterministic trends
    \item Stochastic trends
    \item Trend estimation
    \item Detrending
    \item Variance-stabilizing transformations
\end{itemize}

\subsection{Key Concepts}

A deterministic trend model can be written as

\[
X_t = \beta_0 + \beta_1 t + \epsilon_t,
\]

whereas a stochastic trend may arise from an integrated process such as

\[
X_t = X_{t-1} + \epsilon_t.
\]

A first difference is

\[
\Delta X_t
=
X_t-X_{t-1}.
\]

\subsection{Learning Objectives}

Students should be able to:

\begin{enumerate}
    \item Identify trends in time-series data.
    \item Distinguish deterministic and stochastic trends.
    \item Estimate deterministic trends.
    \item Explain why trends can violate stationarity.
    \item Apply differencing conceptually.
    \item Explain the purpose of variance-stabilizing transformations.
\end{enumerate}

\textbf{Difficulty: High}

\section{Module 4: Linear Processes and ARMA Models}

\subsection{Major Topics}

\begin{itemize}
    \item General linear processes
    \item Moving-average models
    \item Autoregressive models
    \item ARMA models
    \item Stationarity and invertibility
\end{itemize}

\subsection{Moving-Average Model}

An MA($q$) model can be expressed as

\[
X_t
=
\mu
+
\epsilon_t
+
\theta_1\epsilon_{t-1}
+\cdots+
\theta_q\epsilon_{t-q}.
\]

\subsection{Autoregressive Model}

An AR($p$) model has the form

\[
X_t
=
c
+
\phi_1X_{t-1}
+\cdots+
\phi_pX_{t-p}
+
\epsilon_t.
\]

\subsection{ARMA Model}

An ARMA($p,q$) model combines both components:

\[
X_t
=
c+
\sum_{i=1}^{p}\phi_iX_{t-i}
+
\epsilon_t
+
\sum_{j=1}^{q}\theta_j\epsilon_{t-j}.
\]

Using the backshift operator $B$,

\[
\phi(B)X_t
=
\theta(B)\epsilon_t.
\]

\subsection{Learning Objectives}

Students should be able to:

\begin{enumerate}
    \item Construct AR, MA, and ARMA models.
    \item Interpret autoregressive coefficients.
    \item Interpret moving-average coefficients.
    \item Derive model properties from the defining equations.
    \item Understand stationarity conditions for AR processes.
    \item Understand invertibility conditions for MA processes.
    \item Relate model structure to ACF behavior.
\end{enumerate}

\textbf{Difficulty: Very High}

\section{Module 5: Nonstationary Models and ARIMA}

\subsection{Major Topics}

\begin{itemize}
    \item Nonstationarity
    \item Differencing
    \item Integrated processes
    \item ARIMA models
    \item Transformations
\end{itemize}

\subsection{ARIMA Structure}

An ARIMA($p,d,q$) model applies $d$ differences before modeling
the resulting stationary process with an ARMA($p,q$) structure:

\[
\phi(B)(1-B)^dX_t
=
\theta(B)\epsilon_t.
\]

\subsection{Learning Objectives}

Students should be able to:

\begin{enumerate}
    \item Explain why differencing is used.
    \item Determine the appropriate conceptual order of differencing.
    \item Construct an ARIMA model.
    \item Interpret $p$, $d$, and $q$.
    \item Distinguish ARMA and ARIMA models.
    \item Understand the relationship between integration and nonstationarity.
\end{enumerate}

\textbf{Difficulty: Very High}

\section{Module 6: Model Specification and Identification}

\subsection{Major Topics}

\begin{itemize}
    \item Model specification
    \item Model identification
    \item ACF
    \item PACF
    \item EACF
    \item Choosing $p$, $d$, and $q$
\end{itemize}

\subsection{Key Concepts}

The identification process attempts to determine plausible values for
the orders of an ARIMA model:

\[
(p,d,q).
\]

The principal diagnostic tools are:

\begin{itemize}
    \item ACF: Autocorrelation Function
    \item PACF: Partial Autocorrelation Function
    \item EACF: Extended Autocorrelation Function
\end{itemize}

Typical identification logic includes comparing decay and cutoff
patterns in the ACF and PACF.

\subsection{Learning Objectives}

Students should be able to:

\begin{enumerate}
    \item Interpret sample ACF plots.
    \item Interpret sample PACF plots.
    \item Use ACF/PACF behavior to suggest ARMA orders.
    \item Use EACF as an additional identification tool.
    \item Propose candidate ARIMA models.
    \item Recognize that model identification is exploratory rather
          than purely mechanical.
\end{enumerate}

\textbf{Difficulty: Very High}

\section{Module 7: Parameter Estimation}

\subsection{Major Topics}

\begin{itemize}
    \item Yule--Walker estimation
    \item Maximum likelihood estimation
    \item Least squares estimation
    \item Conditional least squares
\end{itemize}

\subsection{Yule--Walker Equations}

For an AR($p$) model, the Yule--Walker equations connect the
autoregressive parameters to the autocorrelation function:

\[
\begin{pmatrix}
1 & \rho(1) & \cdots & \rho(p-1)\\
\rho(1) & 1 & \cdots & \rho(p-2)\\
\vdots & \vdots & \ddots & \vdots\\
\rho(p-1)&\rho(p-2)&\cdots&1
\end{pmatrix}
\begin{pmatrix}
\phi_1\\
\phi_2\\
\vdots\\
\phi_p
\end{pmatrix}
=
\begin{pmatrix}
\rho(1)\\
\rho(2)\\
\vdots\\
\rho(p)
\end{pmatrix}.
\]

\subsection{Learning Objectives}

Students should be able to:

\begin{enumerate}
    \item Estimate AR parameters using Yule--Walker equations.
    \item Explain the principle of maximum likelihood estimation.
    \item Understand least-squares estimation in time-series models.
    \item Distinguish conditional least squares from ordinary regression.
    \item Compare different estimation approaches.
    \item Interpret estimated time-series parameters.
\end{enumerate}

\textbf{Difficulty: Very High}

\section{Module 8: Model Diagnostics and Model Fit}

\subsection{Major Topics}

\begin{itemize}
    \item Residual analysis
    \item Residual autocorrelation
    \item Model adequacy
    \item Goodness of fit
    \item Diagnostic testing
\end{itemize}

\subsection{Core Idea}

After fitting a model,

\[
X_t = \widehat{X}_t + e_t,
\]

the residual sequence $\{e_t\}$ should contain little systematic
time dependence if the model is adequate.

\subsection{Learning Objectives}

Students should be able to:

\begin{enumerate}
    \item Compute and interpret residuals.
    \item Examine residual ACFs.
    \item Determine whether residuals resemble white noise.
    \item Detect remaining serial dependence.
    \item Evaluate model adequacy.
    \item Use diagnostic tests to compare candidate models.
\end{enumerate}

\textbf{Difficulty: High}

\section{Module 9: Forecasting}

\subsection{Major Topics}

\begin{itemize}
    \item Linear forecasting
    \item ARIMA forecasting
    \item One-step-ahead prediction
    \item Multi-step prediction
    \item Prediction intervals
\end{itemize}

\subsection{Forecasting Objective}

Given information through time $t$, construct

\[
\widehat{X}_{t+h|t}
=
E[X_{t+h}\mid X_t,X_{t-1},\ldots].
\]

The forecast error is

\[
e_{t+h}
=
X_{t+h}-\widehat{X}_{t+h|t}.
\]

A prediction interval can be constructed using the forecast-error
variance.

\subsection{Learning Objectives}

Students should be able to:

\begin{enumerate}
    \item Generate one-step-ahead forecasts.
    \item Generate multi-step forecasts.
    \item Forecast using ARIMA models.
    \item Calculate or interpret forecast-error variance.
    \item Construct prediction intervals.
    \item Explain why forecast uncertainty increases with horizon.
\end{enumerate}

\textbf{Difficulty: Very High}

\section{Module 10: Seasonal Box--Jenkins Models}

\subsection{Major Topics}

\begin{itemize}
    \item Seasonal dependence
    \item Seasonal differencing
    \item Seasonal ARIMA
    \item SARIMA forecasting
    \item Seasonal diagnostics
\end{itemize}

\subsection{SARIMA Structure}

A seasonal ARIMA model can be represented as

\[
\Phi(B^s)
\phi(B)
(1-B)^d
(1-B^s)^D
X_t
=
\Theta(B^s)
\theta(B)
\epsilon_t,
\]

where $s$ denotes the seasonal period.

\subsection{Learning Objectives}

Students should be able to:

\begin{enumerate}
    \item Recognize seasonal patterns.
    \item Distinguish ordinary and seasonal differencing.
    \item Explain the components of a SARIMA model.
    \item Specify seasonal AR and MA components.
    \item Forecast using seasonal models.
    \item Diagnose seasonal model adequacy.
\end{enumerate}

\textbf{Difficulty: Very High}

\section{Module 11: Regression Methods for Time Series}

\subsection{Major Topics}

\begin{itemize}
    \item Regression with time-series errors
    \item Intervention analysis
    \item Outlier detection
    \item Spurious correlation
    \item Pre-whitening
\end{itemize}

\subsection{Intervention Analysis}

An intervention model can represent the impact of an external event:

\[
X_t
=
\text{baseline process}
+
\text{intervention effect}
+
\epsilon_t.
\]

\subsection{Spurious Correlation}

Two nonstationary time series may display strong sample correlation
even when their underlying processes are not meaningfully related.

\subsection{Pre-whitening}

Pre-whitening seeks to remove serial dependence from one or more
series before assessing cross-correlation.

\subsection{Learning Objectives}

Students should be able to:

\begin{enumerate}
    \item Explain why ordinary regression assumptions may fail for
          time-series data.
    \item Model interventions.
    \item Identify potential outliers.
    \item Explain spurious correlation.
    \item Explain the purpose of pre-whitening.
    \item Interpret cross-correlation in the presence of serial dependence.
\end{enumerate}

\textbf{Difficulty: Very High}

\section{Module 12: Heteroskedastic Time-Series Models}

\subsection{Major Topics}

\begin{itemize}
    \item Conditional variance
    \item Conditional heteroskedasticity
    \item ARCH processes
    \item GARCH processes
    \item Volatility clustering
\end{itemize}

\subsection{ARCH Model}

An ARCH($q$) model can be represented by

\[
X_t = \sigma_t\epsilon_t,
\]

where

\[
\sigma_t^2
=
\alpha_0
+
\sum_{i=1}^{q}
\alpha_iX_{t-i}^2.
\]

\subsection{GARCH Model}

A GARCH($p,q$) model extends the conditional variance equation:

\[
\sigma_t^2
=
\alpha_0
+
\sum_{i=1}^{q}\alpha_i X_{t-i}^2
+
\sum_{j=1}^{p}\beta_j\sigma_{t-j}^2.
\]

\subsection{Learning Objectives}

Students should be able to:

\begin{enumerate}
    \item Explain conditional heteroskedasticity.
    \item Distinguish constant and time-varying variance.
    \item Explain volatility clustering.
    \item Specify an ARCH model.
    \item Specify a GARCH model.
    \item Interpret conditional variance parameters.
\end{enumerate}

\textbf{Difficulty: Very High}

\section{Module 13: Advanced Topics}

These topics are explicitly identified as topics that may be covered
as time permits.

\subsection{Nonlinear Time Series}

Possible conceptual topics include:

\begin{itemize}
    \item Nonlinear dependence
    \item Nonlinear stochastic processes
    \item Limitations of linear ARMA/ARIMA models
\end{itemize}

\subsection{Spectral/Frequency-Domain Methods}

The frequency-domain perspective represents time-series behavior
in terms of frequencies rather than solely in terms of lags.

Core concepts may include:

\begin{itemize}
    \item Frequency-domain representation
    \item Spectral density
    \item Periodic behavior
    \item Relationship between autocovariance and spectrum
\end{itemize}

\textbf{Difficulty: Very High}

\section{Integrated Model-Building Workflow}

A central practical skill is learning to move from raw data to an
adequate forecasting model.

\subsection*{Step 1: Explore the Data}

\begin{itemize}
    \item Plot the time series.
    \item Examine changes in level.
    \item Examine trend.
    \item Examine seasonality.
    \item Examine changing variance.
\end{itemize}

\subsection*{Step 2: Assess Stationarity}

\begin{itemize}
    \item Examine the mean behavior.
    \item Examine the variance.
    \item Examine serial dependence.
    \item Determine whether transformation or differencing is needed.
\end{itemize}

\subsection*{Step 3: Transform the Series}

Possible operations include

\[
Y_t = g(X_t)
\]

for variance stabilization and

\[
Y_t = \Delta^d X_t
\]

for differencing.

\subsection*{Step 4: Identify Candidate Models}

Use:

\begin{itemize}
    \item ACF
    \item PACF
    \item EACF
    \item Knowledge of the data-generating process
\end{itemize}

to propose candidate ARMA/ARIMA/SARIMA models.

\subsection*{Step 5: Estimate Parameters}

Apply appropriate estimation methods:

\begin{itemize}
    \item Yule--Walker
    \item Least squares
    \item Conditional least squares
    \item Maximum likelihood
\end{itemize}

\subsection*{Step 6: Diagnose the Model}

Analyze residuals:

\[
e_t = X_t-\widehat{X}_t.
\]

Check whether residuals exhibit remaining structure.

\subsection*{Step 7: Refine the Model}

If diagnostics reveal inadequate fit:

\[
\text{Identify}
\rightarrow
\text{Estimate}
\rightarrow
\text{Diagnose}
\rightarrow
\text{Refit}.
\]

This iterative process is the central practical modeling cycle.

\subsection*{Step 8: Forecast}

After selecting an adequate model:

\[
\widehat{X}_{t+h|t}
=
E(X_{t+h}\mid\mathcal{F}_t)
\]

is used for prediction, together with forecast-error uncertainty
and prediction intervals.

\section{Concept Dependency Map}

\begin{center}
\begin{tabular}{>{\raggedright\arraybackslash}p{0.28\textwidth}
                >{\raggedright\arraybackslash}p{0.28\textwidth}
                >{\raggedright\arraybackslash}p{0.28\textwidth}}
\toprule
\textbf{Foundation} & \textbf{Core Model} & \textbf{Application} \\
\midrule
Probability &
Stationarity &
Forecasting \\

Expectation/variance &
AR models &
Seasonal forecasting \\

Covariance &
MA models &
Intervention analysis \\

Correlation &
ARMA models &
Outlier analysis \\

Regression &
ARIMA models &
Pre-whitening \\

Linear algebra &
SARIMA models &
ARCH/GARCH \\

Calculus &
Parameter estimation &
Spectral methods \\
\bottomrule
\end{tabular}
\end{center}

\section{Difficulty Progression}

\begin{longtable}{p{0.08\textwidth}p{0.25\textwidth}p{0.30\textwidth}p{0.20\textwidth}}
\toprule
\textbf{Level} & \textbf{Topic} & \textbf{Main Challenge} & \textbf{Prerequisites} \\
\midrule
1 & Fundamental tools &
Understanding lag-based dependence &
Probability, statistics \\

2 & Stationarity &
Connecting model structure with stationarity &
Probability \\

3 & Trends &
Distinguishing deterministic and stochastic trends &
Regression, stationarity \\

4 & AR/MA &
Understanding recursive and moving-average structures &
Stationarity, linear algebra \\

5 & ARMA &
Combining AR and MA behavior &
AR, MA, stochastic processes \\

6 & ARIMA &
Handling nonstationarity through differencing &
ARMA, trends \\

7 & Identification &
Inferring model structure from ACF/PACF/EACF &
ARMA/ARIMA \\

8 & Estimation &
Understanding time-dependent likelihoods and estimators &
Probability, optimization \\

9 & Diagnostics &
Determining whether a fitted model is adequate &
Identification, estimation \\

10 & Forecasting &
Deriving conditional predictions and uncertainty &
ARIMA, estimation \\

11 & SARIMA &
Combining ordinary and seasonal dynamics &
ARIMA, forecasting \\

12 & Regression methods &
Separating explanatory effects from serial dependence &
Regression, ARIMA \\

13 & ARCH/GARCH &
Modeling conditional rather than constant variance &
Probability, regression, time series \\

14 & Spectral methods &
Understanding the frequency-domain representation &
Autocovariance, Fourier concepts \\
\bottomrule
\end{longtable}

\section{High-Priority Difficult Topics}

The following topics should receive additional study time:

\begin{enumerate}
    \item Weak versus strict stationarity.
    \item Stationarity and invertibility conditions.
    \item Derivation and interpretation of ARMA autocorrelation structures.
    \item Differencing and the relationship between ARIMA and ARMA.
    \item ACF/PACF/EACF-based model identification.
    \item Yule--Walker equations.
    \item Maximum likelihood estimation for dependent observations.
    \item Multi-step ARIMA forecasting.
    \item Forecast-error variance and prediction intervals.
    \item Seasonal differencing and SARIMA specification.
    \item Residual diagnostics and model adequacy.
    \item Spurious correlation in nonstationary series.
    \item Pre-whitening and cross-correlation.
    \item ARCH/GARCH conditional variance models.
    \item Frequency-domain/spectral representations.
\end{enumerate}

\section{Recommended Learning Sequence}

A student studying independently should follow this order:

\[
\boxed{
\text{Prerequisites}
\rightarrow
\text{Fundamentals}
\rightarrow
\text{Stationarity}
}
\]

\[
\boxed{
\text{Stationarity}
\rightarrow
\text{AR/MA}
\rightarrow
\text{ARMA}
\rightarrow
\text{ARIMA}
}
\]

\[
\boxed{
\text{ARIMA}
\rightarrow
\text{Identification}
\rightarrow
\text{Estimation}
\rightarrow
\text{Diagnostics}
}
\]

\[
\boxed{
\text{Diagnostics}
\rightarrow
\text{Forecasting}
\rightarrow
\text{SARIMA}
}
\]

\[
\boxed{
\text{Regression}
\rightarrow
\text{Intervention Analysis}
\rightarrow
\text{Pre-whitening}
}
\]

\[
\boxed{
\text{Conditional Variance}
\rightarrow
\text{ARCH}
\rightarrow
\text{GARCH}
}
\]

\[
\boxed{
\text{Advanced Topics}
\rightarrow
\text{Nonlinear Time Series}
+
\text{Spectral Analysis}
}
\]

\section{Suggested Assessment Alignment}

The syllabus specifies six homework assignments worth 25\% of the
course grade and three midterms worth 25\% each. The second and third
midterms are noncumulative. 

A topic-based assessment structure can therefore be organized as follows:

\begin{longtable}{p{0.18\textwidth}p{0.32\textwidth}p{0.38\textwidth}}
\toprule
\textbf{Assessment} & \textbf{Primary Material} & \textbf{Skills} \\
\midrule
Midterm 1 &
Fundamentals, stationarity, trends, AR/MA/ARMA &
Derivation, calculations, conceptual understanding \\

Midterm 2 &
ARIMA, identification, estimation, diagnostics &
Model construction, identification, estimation \\

Midterm 3 &
Forecasting, SARIMA, regression, ARCH/GARCH &
Integrated model analysis and forecasting \\

Homework &
All modules &
Theory, calculations, modeling, computation \\
\bottomrule
\end{longtable}

\section{Final Competency Profile}

Upon completion of the course, a successful student should be capable
of taking a time series from raw observations through a complete
statistical modeling workflow:

\[
\boxed{
\text{Raw Time Series}
\rightarrow
\text{Exploration}
\rightarrow
\text{Stationarity Assessment}
\rightarrow
\text{Transformation}
}
\]

\[
\boxed{
\rightarrow
\text{Model Identification}
\rightarrow
\text{Parameter Estimation}
\rightarrow
\text{Diagnostics}
}
\]

\[
\boxed{
\rightarrow
\text{Model Selection}
\rightarrow
\text{Forecasting}
\rightarrow
\text{Prediction Intervals}
}
\]

The student should also be able to recognize when standard linear
time-series models are insufficient and identify extensions such as
seasonal models, intervention models, ARCH/GARCH models, nonlinear
models, or frequency-domain methods.

\section{Compact Course Summary}

\begin{center}
\begin{tabular}{p{0.22\textwidth}p{0.65\textwidth}}
\toprule
\textbf{Area} & \textbf{Core Knowledge} \\
\midrule
Foundations &
Mean, variance, covariance, autocovariance, autocorrelation \\

Stationarity &
Strict and weak stationarity, white noise, dependence \\

Trends &
Deterministic/stochastic trends, trend estimation, transformations \\

Linear Models &
Linear processes, AR, MA, ARMA \\

Nonstationarity &
Differencing and ARIMA \\

Identification &
ACF, PACF, EACF, model specification \\

Estimation &
Yule--Walker, maximum likelihood, least squares, conditional LS \\

Diagnostics &
Residual analysis and model-fit assessment \\

Forecasting &
Linear forecasting, ARIMA prediction, prediction intervals \\

Seasonality &
SARIMA, seasonal forecasting and diagnostics \\

Regression &
Intervention analysis, outliers, spurious correlation, pre-whitening \\

Volatility &
ARCH and GARCH \\

Advanced &
Nonlinear time series and spectral/frequency-domain methods \\
\bottomrule
\end{tabular}
\end{center}

\section*{Conclusion}

The course is fundamentally organized around the statistical modeling
cycle

\[
\boxed{
\text{Understand dependence}
\rightarrow
\text{Build a model}
\rightarrow
\text{Estimate it}
\rightarrow
\text{Diagnose it}
\rightarrow
\text{Forecast}
}
\]

with ARMA/ARIMA models serving as the central modeling framework and
SARIMA, regression/intervention, ARCH/GARCH, nonlinear, and spectral
methods providing important extensions.

\end{document}